\documentclass[11pt]{article}

% Packages
\usepackage[a4paper, margin=1in]{geometry} % Adjust margins
\usepackage{setspace} % For line spacing
\usepackage{tocloft} % Optional: better spacing for TOC
\usepackage[backend=biber,style=apa]{biblatex} % Load biblatex with APA style
\addbibresource{references.bib} % Link to your .bib file

% Set line spacing
\onehalfspacing % or use \doublespacing for double spacing

% Title information
\title{Preferences on Tax Enforcement and Tax Administration}
\author{Tim Kircali, Stefanie Stantcheva and Matthias Weber}
\date{\today} % Use \date{} for no date

\begin{document}

\maketitle

\section{Introduction}
Tax revenue is essential for modern governments to provide their functions, reaching from the provision of safety over infrastructure, and education to healthcare and social security. Because it is of crucial importance that the taxes that are owed are also paid, governments and tax administrations have to decide how to design an effective tax system to enforce tax revenue. Besides choosing the level of tax enforcement\footnote{Existing evidence suggests that additional investments in tax enforcement are welfare improving (\cite{boning2025welfare})}, governments have to decide which instruments tax authorities can and should use to enforce taxes. For example several countries implemented automatic exchange of information in tax systems allowing tax administrations to easily validate information while other countries tax authorities don't have access and mostly have to rely on trust to the taxpayer filing his returns correctly. And if we switch the viewpoint to a service perspective of a tax administration, tax authorities have to decide which services they want to provide for tax filers. While some countries offer prepopulated tax returns and free government provided tax software, other countries' taxpayers face considerable tax compliance costs and have to pay for tax software by themselves.%\footnote{e.g. American taxpayer has 12 hours + 270\$ tax filing costs} 

One possible approach to understand and decide how a tax system should be designed and which services it should offer, is to directly elicit a society’s preference for it. In a democracy, one may argue that the level \& design of tax enforcement and which services a tax administration should offer, should correspond to a population’s preference. Similar to previous studies analyzing people's policy support using surveys (\textcite{kuziemko2015elastic} for the demand for redistribution, \textcite{stantcheva2021understanding} for the preference on the tax schedule), we analyze how people's preferences for the design of tax enforcement and government-provided tax filing services look like. In a large survey on citizens from the US, and possibly other countries, we ask people about their views on tax enforcement policies. Should governments expand automatic exchange of information? And if we switch to the tax administration's service perspective, should tax administrations offer (free) tax software and prepopulated tax returns to reduce tax filing costs? These questions are closely related, because for a government to be able to offer pre-populated tax returns, it must have automatic access to tax-related information.

And how do people reason about these policy views? What are the considerations driving policy support and what is the role of (mis)information? There might be a big number of considerations underlying support for tax enforcement policies, however, we try to focus on 4 convincing primary factors underlying policy support: Data privacy concerns, considerations on tax evasion and tax revenue, considerations regarding fairness \& inequality and tax filing concerns. On top of these underlying considerations, (mis)information and political leanings and government views can play a strong role in policy preference formation. To shed more light on how these underlying considerations drive policy views we include questions to elicit them. \footnote{And by collecting additional background information on respondents, we can analyse how considerations differ by the socioeconomic and political characteristics of respondents.} Finally, to understand the role of information and to establish a causal link between considerations and policy views, we experimentally provide respondents with information specifically addressing these 4 certain underlying considerations. Doing that we can observe how considerations of treated groups change and more interestingly how this translates into changing policy views. These changes help us to understand which of the underlying considerations are primarily driving policy views. \\[0.5em]
\textbf{Related Literature:} First, this project is related to the broad literature on measuring people's support for policies using experimental surveys. While the literature measures policy preferences in several areas, \textcite{stantcheva2021understanding}, who elicits preferences on tax policy in general, is most closely linked. While she concentrates on the question which tax schedule of income and wealth taxes people prefer, we try to understand people's views on which instruments tax administration should use to enforce that tax schedule. 
Second, this project relates to the literature on tax administration, tax enforcement and tax evasion, and especially, on the literature on the role of automatic exchange of information and third-party reporting in reducing tax evasion. While \textcite{kleven2011unwilling}, \textcite{pomeranz2015no}, \textcite{slemrod2017does} show that automatic exchange of information and third-party reporting is effective in reducing tax evasion, our research tries to uncover why many countries' parliaments and societies oppose further expanding automatic exchange of information into their tax systems. 
Finally, our project contributes to the literature on tax filing costs and the governments role in providing services to reduce them. \textcite{benzarti2020taxing} observes how taxpayers choose between itemizing deductions and claiming the standard deduction and finds that taxpayers forgo large tax savings to avoid compliance costs. \textcite{hauck2024optional} show that in Germany non-filing of tax declaration is common, in particular at the lower end of the income distribution. Therefore non-filing reduces the effectiveness of redistribution. \textcite{benzarti2024rising} use word counts of the tax codes in several countries dating as far back as the early 1980s as a proxy for tax-compliance costs and show that tax compliance costs have been steadily increasing. Using a survey on US taxpayers they show that the majority would be willing to pay for simplifying the tax system and adopting prepopulated tax returns. Our survey strongly expands this research with several aspects. First it also looks at peoples' demand for tax software. And second it elicits people's views about high automatic exchange of information, which is in itself a prerequisite for government provided prepopulated tax returns. Finally, by experimentally providing information we can understand the role of (mis)information in shaping policy decisions.  \\[-0.5em] 

In the rest of the paper we will first discuss the data sources and the survey design. Then we will discuss expected results before finally concluding. 
\section{Conceptual Framework:} How do people reason about political views? What are the considerations that drive political support for automatic exchange of information and their demand for governments to provide tax software and pre-populated tax returns? From an economic perspective, there are several considerations that could underlie support for tax enforcement policies. However, the most closely related ones would be the following. \\[0.5em]
\textbf{Data Privacy:} First, people may generally dislike (government) institutions collecting and having access to private (financial) data. If the AEI policy in question either increases the amount of data exchanged with (government) institutions or increases the amount of data exchanged between institutions, people might oppose AEI policies due to financial privacy considerations. \\[0.5em]
\textbf{Evasion and Revenue:} Second, people might be interested in governments effectively collecting tax revenue because they believe that government tax revenue is important. Existing evidence shows that audit rates are low, tax evasion is high, and that automatic exchange of information generally helps tax administrations to reduce tax evasion (Kleven et al. 2011, Pomeranz 2015, Slemrod et al. 2017). These are strong arguments in favour of stronger AEI policies. \\[0.5em]
\textbf{Fairness and Inequality:} Third, and closely related, people might be interested in everyone paying their fair share of taxes. The fact that tax evasion is much more pronounced at the top of the income distribution, and the consideration that honest taxpayers have to finance the tax gaps created by evaders, would be reasons to support AEI policies. \\[0.5em]
\textbf{Tax Filing:} Fourth, people may have considerations about tax filing and compliance costs that influence their political views. Indeed, \textcite{benzarti2024rising} show that US taxpayers perceive that tax complexity and filing costs have increased and that the majority would be willing to pay for a simplification of the tax system and the introduction of prepopulated tax returns. People may also be concerned about making mistakes when filing their taxes. The fact that automatic exchange of information would help tax administrations to offer services such as free government tax software and prepopulated tax returns, which could significantly reduce tax filing costs, would be a reason for AEI policy support. \\[0.5em]
\textbf{Knowledge:} Knowledge may not be directly related to people's preferred AEI policy, but it might shape policy views through the considerations discussed above. Consider two equivalent people with a strong preference for prepopulated tax returns and a weaker preference for privacy, where the only difference between them is that person A understands the importance of AEI for governments to provide pre-populated tax returns, while person B does not. In this scenario, person A might support the AEI policy because of his demand for pre-populated tax returns, while person B might not support it because of his lack of understanding the link between AEI and prepopulated tax returns. \\[-0.5em]

Using an experimental survey approach, we will measure underlying considerations and how they are related to policy views. To understand the role of information, we elicit people's prior knowledge about the economic factors underlying AEI policies and analyse how knowledge levels are related to considerations and policy views. To establish a causal link between considerations and policy views, we experimentally provide respondents with information that specifically addresses one of the underlying considerations mentioned above.  This allows us to observe how the considerations of the treated groups change and, more interestingly, how this translates into changes in policy views. These changes help us to understand which of the underlying considerations are primarily driving policy views. 
\section{Data and Survey Design}
\subsection{Data Collection}
The core data will come from a survey conducted in end 2025 on U.S. residents aged 18 to 69. We consider expanding the survey to other countries such as Germany, Switzerland and Denmark. The planned sample sizes lie between 2500 (for the U.S. Sample) and 6000 (for all 4 countries). We will use a commercial survey provider such as Bilendi, Dynata or YouGov. These companies have large pools of vetted survey respondents that come from different panels, who can freely log in and take surveys in exchange for rewards. Typically they are consumer surveys, that are, related to consumer products or experiences. On the entry page we just give them information about the length of the survey to avoid selection based on the topic. Then they are guided trough a consent page and after that trough some screening questions that ensure that the final sample is representative along gender, age, and education/income. 
\subsection{Survey Structure} 
\textbf{Demographic Background and Political Leanings:}
To begin the survey, we ask people about their demographic background, about their political orientation and their views about the government. Due to its importance, political leaning is elicited with more than one question: For example we ask where they see themselves on a left/right spectrum and how they voted in the last elections. To understand their views about the government, we ask them if they trust the government, what the scope of government intervention should be and ask about their perceptions of governments capabilities providing consumer friendly administrative solutions to reduce compliance costs of citizen. The last question aims to elicit if people think that the government would be capable of providing user friendly tax software or prepopulated tax returns without mentioning these words to avoid priming. \\[0.5em]
\textbf{Knowledge}
These questions are about factual knowledge that are relevant to understand the AEI policies in question. We ask respondents if they know what automatic exchange of information with tax administration is and how it works. We ask them about their perceptions of the tax gap and where along the distribution it is being evaded. We consider using incentivised questions for a subsample we to analyse how precision of people's responses change.  \\[0.5em]
\textbf{Information Treatments}
Respondents are randomly split into 6 groups. Each group of respondents will see either one of 4 different instructional information video treatments, a full info treatment including all information from the treatments before, or no information at all (control). These information treatments each address certain underlying considerations. More details about the treatments will be given at a later stage.\\[0.5em]
\textbf{Considerations on AEI Policies}
In this section we ask people about their considerations and expectations about the economic impact of very general AEI policies. For example we refer to an expansion of automatic exchange of information and describe it as "more tax-relevant data is being exchanged between institutions and tax administration". First we ask about their data privacy concerns on these policies. Second, we ask about their expectations of the economic impact on factors such as tax evasion and tax revenue. Third, we ask about their concerns regarding fairness and inequality. We ask them if they think income inequality is concerning and if they think it can be reduced with automatic exchange of information. And we ask them if they think the current tax system is fair and if AEI helps to make it fairer. Finally, we ask about their tax filing and tax compliance considerations. For example we ask them about their concerns about tax filing costs and if they think that fearing to make mistakes in tax filing is a legitimate concern.  And we ask them if they think that automatic exchange of information might help to reduce tax filing costs and reduce the risk of making mistakes. \\[0.5em]
\textbf{Support for AEI Policies}
In the last block we ask them about their preferences on policies related to automatic exchange of information. Specifically we ask them if they would support an expansion of automatic exchange of information. To further understand reasons for their support we ask them if they would support expansion of AEI conditional on certain elements, such as if the policy concentrates on high net worth individuals or if it includes corporate tax payers. Finally, we move on to understand their support for government efforts to provide tax software and prepopulated tax returns, where we ask them how much they would value free public tax software and prepopulated tax returns. \\[0.5em]
\textbf{Final Questions}
At the end of the survey we ask some finalizing questions such as if respondents think the survey was biased (left/right) or if they would be willing to pay for information. 
\subsection{Information Treatments}
We use four information treatments in form of illustrative video treatments that are very closely related to one of the primary considerations of policy support. We use this strategy to understand how specific information on certain considerations translates into policy views. Additionally, one group will see a full information treatment which includes the information of all  \\[0.5em]
\textbf{Data Privacy Treatment:}
In the first information treatment we illustrate tax relevant players and institutions and how data is being exchanged between them. Then we explain how information exchange would change if automatic exchange of information is expanded. For that we might use comparisons to other countries and explain how it works there. The primary goal of the data privacy treatment is to shock people’s knowledge and perceptions about automatic exchange of information, specifically addressing their data privacy considerations and distinguishing them from other economic considerations. \\[0.5em]%The underlying assumption for that treatment is that people might generally not understand what automatic exchange of information is and how it works. They might not be aware of the fact that several (governmental) institutions already have data on them and that they are just not being automatically exchanged. They might not know that tax authorities could still get access to data under serious suspicion and that honest tax payers end up giving them the data anyway.  
\textbf{Tax Filing Treatment:}
The second treatment discusses concerns about tax filing costs such as the fact that an average American taxpayer has 12 hours + 270 \$ tax filing costs. It explains how (guided) tax software works and what prepopulated tax returns looks like. It emphasizes that tax software and prepopulated tax returns help tax filers to make fewer mistakes and significantly reduces their tax filing costs. And importantly, we highlight the importance of AEI to the tax administration in providing prepopulated tax returns. The primary goal of the tax filing treatment is to shock people's knowledge and perceptions about tax filing costs and tax filing services and how they are related to the existence of automatic exchange of information. \\[0.5em]%And an important goal would be to isolate this effect from other economic considerations. The underlying assumption of this treatment is, that people might not know about government provided tax software or tax software in general and might not know about the existence of prepopulated tax returns in other countries. And importantly, they might not understand the importance of automatic exchange of information for such government provided services.  
\textbf{Tax Evasion Treatment:}
The third treatment provides respondents with information about the economic effects of automatic exchange of information in tax enforcement. We give them information about the size of the tax gap, about the limited resources of IRS and that audit levels are generally low. We explain that tax evasion is rare in areas which are covered with AEI and point out that AEI is therefore directly effective in reducing tax evasion. The primary goal of the tax evasion treatment is to shock people’s knowledge and perceptions about the economic effects of tax enforcement and automatic exchange of information without affecting considerations related to fairness \& inequality. \\[0.5em]% The underlying assumption of this treatment is, that people might have misperceptions about the extend of tax evasion,  probably don’t know that the audit levels are very low, and finally might not know that AEI is very effective in reducing tax evasion. 
\textbf{Inequality Treatment}
In the last treatment we again give them information about tax evasion but explain the fact the tax gap is much more pronounced for high-income individuals and reasons why this is the case. We explain fairness considerations of tax evasion such as the fact that honest taxpayers have to finance the fiscal gap caused by tax evaders. The primary goal of the inequality treatment is to shock people’s knowledge and perceptions about fairness and distributional considerations of tax enforcement policy and automatic exchange of information. 
\section{Potential Results and Concluding Remarks}
The results of this experimental survey design will help us understand whether people support or oppose automatic information exchange, and what considerations drive these policy views. For example, an interesting finding might be that privacy considerations are strongly pronounced and correlate with opposition to AEI policies. However, if respondents are provided with information about other considerations, these may contribute much more strongly than data privacy considerations to changes in policy views. 

For example, people may have strong concerns about the cost of filing taxes, but prior to the provision of information, they may not understand the importance of AEI for governments to provide services such as pre-populated tax returns. Alternatively, they might have a strong preference for governments to collect tax effectively, but strongly misperceive the size of the tax gap. If they were provided with information about the true extent of tax evasion, their policy preferences might change significantly. 

Understanding their underlying reasoning and the misperceptions involved in forming policy views is crucial to understanding what policies a government should propose. It helps policy makers to decide what elements an AEI policy should include and what elements the public needs to be informed about in order to gain support for the policy. Ultimately, even if the policy debate about automatic exchange of information is mostly tied to people's data privacy concerns, it could be the case that informed people might rather connect the topic to tax filing considerations.

\newpage
\printbibliography
\end{document}